% Generated draft: requires mathematical and source review.
\documentclass[11pt]{article}
\usepackage[margin=25mm]{geometry}
\usepackage{fontspec}
\setmainfont{Latin Modern Roman}
\usepackage{amsmath,amssymb,mathtools,unicode-math}
\setmathfont{Latin Modern Math}
\usepackage{enumitem,needspace,xurl,hyperref}
\setlength{\parindent}{0pt}
\setlength{\parskip}{5pt}
\setlength{\emergencystretch}{3em}
\newcommand{\N}{\mathbb{N}}
\newcommand{\Z}{\mathbb{Z}}
\newcommand{\Q}{\mathbb{Q}}
\newcommand{\R}{\mathbb{R}}
\newcommand{\C}{\mathbb{C}}
\newcommand{\F}{\mathbb{F}}
\newcommand{\A}{\mathbb{A}}
\newcommand{\PP}{\mathbb{P}}
\newcommand{\E}{\mathbb{E}}
\newcommand{\eps}{\varepsilon}
\newcommand{\dd}{\,\mathrm{d}}
\begin{document}
\section*{UnsolvedMath ID: 30003337}
\textbf{Source problem code:} catalogue entry
\subsection{Definitions and mathematical statement}
Let \(G\) have minimum degree \(D\), and form a random graph \(H\) by retaining each edge of \(G\) independently with probability \(1/2\), keeping the vertex set unchanged. Under the conventional interpretation that \(G\) is a finite, nonempty, simple undirected graph, define
\[
\delta(G)=\min_{v\in V(G)}\deg_G(v)=D,
\qquad
d(H)=\max_{\substack{F\subseteq H\\V(F)\ne\varnothing}}\delta(F).
\]
Here \(F\) ranges over subgraphs with at least one vertex, and \(d(H)\) is a random variable. The source does not explicitly specify finiteness, simplicity, or directedness, so this graph convention requires confirmation.

The supplied problem asks how large \(d(H)\) must be in terms of \(D\), uniformly over the intended class of graphs \(G\). In particular, it asks whether
\[
d(H)\geq \frac{D}{3}-o(D).
\]
The asymptotic notation would mean an error \(r(D)\) satisfying \(r(D)/D\to0\) as \(D\to\infty\). However, the source does not specify the probabilistic meaning of this inequality, or whether the error is uniform over \(G\). For example, an assertion about \(\mathbb E[d(H)]\) and an assertion that \(d(H)\) exceeds a threshold with probability tending to one are different statements. No such interpretation is selected here. Consequently, the displayed inequality records the supplied target but is not a fully quantified conjecture until these qualifications are resolved.
\subsection{Short English statement}
Independently keep half the edges of a graph of minimum degree \(D\), and consider the greatest minimum degree among its nonempty subgraphs. The question asks for a universal lower bound, especially \(D/3-o(D)\). The intended probability guarantee, graph class, and uniformity of the error term are unspecified.
\subsection{Sources}
\begin{itemize}
\item Local UnsolvedMath registry, supplied source record for permanent ID \(30003337\), titled \emph{Minimum Degree after Random Edge Percolation}; exact source problem code \texttt{catalogue entry}. The record belongs to the Oberwolfach Reports Open Problems collection and supplies proposed year \(2017\), without a named proposer.
\item Supplied report citation: \emph{Combinatorics (2018)}, Oberwolfach Reports, DOI \texttt{10.4171/owr/2017/1}, pp.\ 5--81. Original OWR report: \url{https://doi.org/10.4171/owr/2017/1}. The supplied title year is retained despite its apparent mismatch with the DOI year and proposed year.
\item Supplied literature reference: Bukh--Krivelevich--Narayanan, ``colouring random subgraphs'': \url{https://arxiv.org/abs/2312.08340}.
\item Literature assessment dated exactly \texttt{2026-08-21}: the supplied curation labels the problem open and states that the universal lower bound remains \(D/4\) minus lower-order terms, while a 2025 construction gives an upper bound \(D/3\) plus lower-order terms for infinitely many degrees; it describes the exact asymptotic value as open. This is supplied dated curation, not an independently verified literature assessment. Its probability guarantees and other technical hypotheses are not supplied.
\end{itemize}
\end{document}
